% ===========================================================================
% Category: Recurrent Neural Network
% Generated from ML_Lab_Manual_Completed.md - edit the Markdown and re-run
% scripts/build_latex_manual.py instead of editing this file by hand.
% ===========================================================================

\chapter{Experiment 14: Recurrent Neural Network
(RNN)}\label{experiment-14-recurrent-neural-network-rnn}

\textbf{Category:} Sequence modeling

\section{1. Objective}\label{objective}

Build recurrent networks for sequential data: (a) sequence
classification and (b) time-series forecasting, comparing vanilla RNN
with LSTM.

\section{2. Datasets Used}\label{datasets-used}

\begin{enumerate}
\def\labelenumi{(\alph{enumi})}
\tightlist
\item
  Synthetic sequence benchmark: 1,200 sequences of length 20, label = 1
  if the sum is positive (1,000 train / 200 test). (b) Monthly Airline
  Passengers (\passthrough{\lstinline!data/airline\_passengers.csv!}):
  144 months, 12-month sliding windows, chronological 80/20 split.
\end{enumerate}

\section{3. Required Libraries}\label{required-libraries}

PyTorch, numpy, pandas, matplotlib (the manual's TensorFlow/Keras
example is implemented equivalently in PyTorch, the framework installed
in this environment)

\section{4. Brief Theory}\label{brief-theory}

An RNN maintains a hidden state h\_t = tanh(W x\_t + U h\_\{t−1\} + b)
so the output depends on the sequence history; training uses
backpropagation through time. Vanilla RNNs suffer from vanishing
gradients over long sequences; LSTM gates (forget/input/output +
additive cell state) preserve long-range information.

\section{5. Procedure}\label{procedure}

\begin{enumerate}
\def\labelenumi{\arabic{enumi}.}
\tightlist
\item
  Generate the synthetic sequence classification data and split
  1,000/200.
\item
  Build RNN and LSTM classifiers (hidden 32, linear sigmoid output);
  train 10 epochs with Adam and BCEWithLogitsLoss.
\item
  Report test accuracy for both.
\item
  For forecasting: normalize the passenger series, build 12-month
  windows, split chronologically.
\item
  Train RNN and LSTM forecasters (hidden 64, 2 layers, MSE, Adam, 120
  epochs).
\item
  Report RMSE/MAE and plot the forecasts.
\end{enumerate}

\section{6. Reference Implementation}\label{reference-implementation}

```python \#\# 10. Complete Code Listing

\emph{Full runnable program for this experiment, extracted from
\passthrough{\lstinline!notebooks/06\_recurrent\_neural\_network.ipynb!}
(executed; results above are its actual output).}

\textbf{Listing 14.1 --- Core numerical and plotting libraries}

```python
